\subsection{An obstruction already in the first filtered differential}

\begin{theorem}[A-024: \status{Refuted}; relative numerical first-stage data]
\label{sa08:A024}
Put $R=\operatorname{Num}(\Z T)$, with augmentation
$\epsilon:R\to\Z$ given by $T\mapsto0$, and put $I=\ker\epsilon$.
Let $\Phi_n=\Z^{n+1}$ be the standard simplicial module with marked
vector $T$, and let $\Psi=\Phi/\Z T$.
There do not exist the following data:
\begin{enumerate}
\item The simplicial $R$-algebra
$B^0=\operatorname{Num}(\Phi)$, with its usual polynomial-degree
filtration and its constant relative coefficient algebra $R$.
\item A simplicial $R$-module $B^1$ with an exhaustive increasing
filtration by simplicial abelian subgroups $F_wB^1$, indexed by
$w\ge0$, such that $F_0B^1$ is simplicially acyclic and
\[
 \operatorname{gr}_wB^1\cong
 \Gamma^{w-1}_{\Z}(\Phi)\otimes\Psi
 \quad(w\ge1)
\]
as simplicial abelian groups.
\item A unital simplicial $R$-bilinear action
$\mu:B^0\otimes_R B^1\to B^1$, and a filtration-preserving
simplicial $R$-linear map $d:B^0\to B^1$ which vanishes on $R$
and whose weight-one associated-graded map is the quotient
$\Phi\to\Psi$.
\end{enumerate}
No associativity of $\mu$, Leibniz rule for $d$, or prescribed
associated-graded product is needed for this obstruction. In
particular it applies to a strict relative filtered DGA, or to a
strict operadic $E_\infty$ model, with these first two rows.
\end{theorem}

\begin{proof}
Write $N_s$ for homological simplicial normalization. The reduced
numerical-polynomial functor
$\operatorname{Num}^+(V)=\ker(\operatorname{Num}(V)\to\Z)$
sends zero to zero. Applying it to the additive simplicial
contraction of $\Phi$ gives a simplicial contraction of
$\operatorname{Num}^+(\Phi)$. Thus
\[
 H_i(N_sB^0)=
 \begin{cases}\Z,&i=0,\\0,&i>0,\end{cases}
 \qquad (N_sB^0)_0=R,
 \qquad \operatorname{im}(\partial:(N_sB^0)_1\to R)=I.
\]
Here the map on $H_0$ is evaluation at zero, and $R$ acts on that
$\Z$ by $\epsilon$. The integral linear coordinates
$T,x_1,\ldots,x_n$ give
$B_n^0=R\otimes_{\Z}\operatorname{Num}(\Z^n)$, which is free
as an $R$-module. Therefore $N_sB^0$ is a nonnegative projective
$R$-resolution of $\Z$.

For $w>1$, the simplicial module
$\Gamma^{w-1}(\Phi)\otimes\Psi$ is acyclic: the positive divided
power of $\Phi$ has an additive simplicial contraction, and its
tensor product with $\Psi$ retains that contraction. The weight-one
module $\Psi$ has normalization $\Z[1]$. The exact sequences for
successive finite filtration stages, followed by exactness of
filtered colimits, consequently give
\[
 H_i(N_sB^1)=
 \begin{cases}\Z,&i=1,\\0,&i\ne1.\end{cases}
 \tag{A-24a}
\]
Moreover $H_1(N_sF_1B^1)\to H_1(N_sB^1)$ and
$H_1(N_sF_1B^1)\to H_1(N_s\operatorname{gr}_1B^1)$ are
isomorphisms.

We next prove that $I$ annihilates $H_1(N_sB^1)$.
Let $f\in I$ and choose $u\in(N_sB^0)_1$ with $\partial u=f$.
Let $z\in(N_sB^1)_1$ be a cycle. The explicit shuffle expression
\[
 v=\mu(s_1u,s_0z)-\mu(s_0u,s_1z)\in B_2^1
\]
has $d_1v=d_2v=0$, by the simplicial identities and the normalized
conditions $d_1u=d_1z=0$. Its remaining face is
\[
 d_0v=\mu(s_0f,z)=fz,
\]
since $d_0z=0$ and the action extends the relative scalar action.
Thus $v\in(N_sB^1)_2$ and $fz$ is a boundary. This argument uses
only the unital bilinear action specified in the theorem.

For $f\in I$, define
\[
 h(f)=[du]\in H_1(N_sB^1),\qquad \partial u=f.
 \tag{A-24b}
\]
This is defined because $\partial(du)=d(f)=0$. If two choices of
$u$ have the same boundary, their difference is a cycle in
$(N_sB^0)_1$, hence is a boundary; applying $d$ shows that their
images define the same class. The map $h:I\to H_1(N_sB^1)$ is
therefore well defined and $R$-linear. Since $I$ annihilates its
target, it vanishes on $I^2$.

In simplicial degree one choose the linear coordinate $u=X_1$
with $d_1u=0$ and $d_0u=T$. The leading-symbol hypothesis says that
$du$ maps to the normalized generator of $\Psi_1$ in
$\operatorname{gr}_1B^1$. By the isomorphisms following (A-24a),
we may orient its homology so that
\[
 h(T)=[du]=1\in\Z.
\]
The numerical polynomial $\beta_2(T)=\binom T2$ belongs to $I$
and satisfies the exact identity
\[
 2\beta_2(T)=T^2-T.
\]
Applying $h$, and using $h(T^2)=0$, gives
$2h(\beta_2(T))=-1$ in $\Z$, a contradiction.
\end{proof}

\paragraph{Derived formulation and its precise scope.}
The same obstruction has a form that does not require a strict
representative of the action. In the derived category $D(R)$, use
homological notation for shifts, so that $\Z[1]$ has homology
$\Z$ in degree one. Then
\[
 \operatorname{Hom}_{D(R)}(\Z,\Z[1])
   =\operatorname{Ext}^1_R(\Z,\Z)=0.
 \tag{A-24c}
\]
Indeed the exact sequence $0\to I\to R\to\Z\to0$ identifies
this Ext group with $\operatorname{Hom}_R(I,\Z)$: restriction
$\operatorname{Hom}_R(R,\Z)\to\operatorname{Hom}_R(I,\Z)$ is
zero because $I$ acts by zero on $\Z$. To see that the latter Hom
group vanishes, let $h:I\to\Z$ be $R$-linear. The numerical identity
\[
 n\beta_n(T)=(T-n+1)\beta_{n-1}(T)
\]
gives, for $n\ge2$,
\[
 n h(\beta_n)=-(n-1)h(\beta_{n-1})
             =(-1)^{n-1}h(T).
\]
Hence $h(T)$ is divisible by every positive integer, so is zero;
then each $h(\beta_n)$ is zero. The $\beta_n$, $n\ge1$, form
an integral basis of $I$, proving (A-24c).

Equivalently, for $P\simeq\Z$ and $Q\simeq\Z[1]$ in $D(R)$,
where the $R$-action on the homology of $Q$ is the augmentation,
there is no morphism $P\to Q$ with nonzero $T$-transgression.
Here the transgression is the map
$I\to\Z$ corresponding to its Ext class through the preceding
exact sequence; on a projective resolution with the relative
coefficient map null, it is exactly (A-24b).
A homotopy-coherent unital action of the simplicial algebra $B^0$
on $B^1$ also forces the augmentation action on its homology:
after normalization $B^0\simeq\Z$ as a derived $R$-algebra.
Thus coherent module operations cannot remove the contradiction
while preserving this nonzero transgression.

This derived statement has an explicit limit. An abstract filtered
$E_\infty$ algebra supplied only with an associated-graded
\emph{total derived object} need not specify separate
exterior-degree-zero and exterior-degree-one simplicial modules,
or the above first differential and transgression. Statement
\ref{sa08:A024} excludes enhanced models that retain this specified
relative first-stage data. It does not silently identify every
abstract filtered totalization with such data.

The finite certificate
\texttt{code/D\_first\_stage\_obstruction.py} verifies the numerical
identities used in the calculation and exhausts the two possible
residues of $h(\beta_2)$ modulo two. Neither solves
$2h(\beta_2)=-1$. The all-index argument establishing (A-24c) is
proved above and is not inferred from that finite test.
